\documentclass[11pt]{article}

\input{0-parameters.tex}

\begin{document}


\section{Empirical validation of the framework}

In this appendix, we \textbf{realistically quantify the errors introduced by  assumptions} to forecast industry-level productivity growth and Domar aggregation, using back-testing, to provide prediction intervals. We are comparing \textbf{model forecasts} with \textbf{actual observed outcomes}.


\subsection{Empirical Data}
\label{sect_empirical_data}

\subsubsection{Set-up} 

We compute capital values at current prices, factor shares, and total factor productivity (TFP) levels. Let $K$ denote capital, $L$ labour, $Y$ value added, $X$ intermediate inputs, $Q$ output, and $A$ total factor productivity. Let 
$i$ index industries and $c$ index countries. For any variable $W \in \{K, L, Y, X, Q\}$, the symbol $W_{ic,t}$ denotes its quantity in sector $i$ and country $c$ at time $t$. The associated price is denoted by $p_{W,ic,t}$, and the nominal value by $V_{W,ic,t}$, such that:
\[
V_{W,ic,t} = p_{W,ic,t} \cdot W_{ic,t}.
\]

\noindent

To test our equations for Domar weights and the productivity $\zeta$, we rely on real and nominal growth rates. We use country-specific proxies for the GDP deflator in each dataset to convert nominal values into real terms. While these proxies broadly capture economy-wide price trends, they may not fully reflect sectoral heterogeneity—particularly in non-market services and smaller industries outside the manufacturing and goods-producing sectors.
\ER{Say more precisely where.}

\noindent
We compute the nominal value of capital, $V_{K,ic,t}$, as the residual between value added and labour costs
\begin{equation}
    V_{K,ic,t} = V_{Y,ic,t} - V_{L,ic,t}.
\end{equation}


\noindent
Input cost shares are denoted by $\alpha$ and computed as
\begin{equation}
    \alpha_{L,ic}(t) = \frac{V_{L,ic,t}}{V_{Y,ic,t}}, \quad
    \alpha_{K,ic}(t) = \frac{V_{K,ic,t}}{V_{Y,ic,t}}, \quad
    \alpha_{X,ic}(t) = \frac{V_{X,ic,t}}{V_{Y,ic,t}}.
\end{equation}

\noindent
Total factor productivity (TFP) growth is computed using a translog index based on input growth rates and their average cost shares:
\begin{equation}
\label{eq_tfp_translog}
\begin{aligned}
    \hat{A}_{ic,t} =
     \hat{Q}_{ic,t} -
     \hat{L}_{ic,t} \cdot \frac{\alpha_{L,ic,t} + \alpha_{L,ic,t-1}}{2} -&
    \hat{K}_{ic,t} \cdot \frac{\alpha_{K,ic,t} + \alpha_{K,ic,t-1}}{2} \\&-
     \hat{X}_{ic,t} \cdot \frac{\alpha_{X,ic,t} + \alpha_{X,ic,t-1}}{2},
\end{aligned}
\end{equation}
\noindent
where $\hat{W}_{ic,t}$ denotes the log growth rate of input $W$ in industry $i$, country $c$, and year $t$.

\noindent
TFP levels $A_{ic,t}$ are then recovered recursively, normalised to 100 in the initial period:
\begin{equation}
    A_{ic,t} =
    \begin{cases}
        100, & \text{if } t=0, \\
        \exp\left( \ln A_{ic,0} + \sum_{t'=1}^{t} \ln \hat{A}_{ic,t'} \right), & \text{if } t > 0.
    \end{cases}
\end{equation}

\noindent
This allows computation of long-run TFP growth between $t$ and $t{+}\tau$ as:
\begin{equation}
    \hat{A}_{ic,t,\tau} = \ln A_{ic,t+\tau} - \ln A_{ic,t}.
\end{equation}

\noindent
We compute the Domar weight of industry $i$ in country $c$ at time $t$ as the ratio of the nominal value of the output $V_{X,ic,t}$  and the country GDP,$V_{Y,c,t}$, that is the aggregated value added in country $c$,
\begin{equation}
    \lambda_{ic,t} = \frac{V_{Q,ic,t}}{V_{Y,c,t}}= \frac{p_{Q,ic,t}Q_{ic,t}}{p_{Y,ic,t}Y_{ic,t}},
    \label{eq:lambdajct}
\end{equation}
\noindent
where $V_{X,ic,t}$ is the nominal output of industry $i$ in country $c$, and $V_{Y,c,t}$ is aggregate value added in country $c$.

\subsubsection{KLEMS}

The \textit{KLEMS} database provides data for 29 countries and 56 sectors over the period 1996–2020, with most observations concentrated between 2009 and 2018. It covers time spans of 10 to 25 years depending on the country. We use the ``Growth Accounts Basic'' and the ``National Accounts'' files of the ``Statistical Module''. 
Data quality varies considerably across countries and sectors. For example, the water transport sector in Slovakia records total factor productivity (TFP) growth of 139\% in 2007, while professional services in Ireland show TFP growth of 97\% in 2015. Following visual inspection, we exclude Slovakia, Greece, and Ireland due to implausible values. We also exclude the aggregated Union entries (covering 11 or 12 countries), as they duplicate information from individual countries. We restrict the analysis to data up to 2019, as the COVID-19 pandemic introduced substantial negative productivity shocks in 2020.
We use the value added price index (VA\_PI) of the aggregate sector “TOT” as an approximation of the economy-wide GDP deflator.

The KLEMS data codes correspond to our variables as follows
\begin{align*}
&V_Q = \text{GO\_CP} \hspace{1cm}& 
&Q = \text{GO\_Q} \hspace{1cm}& 
&P_Q = \text{GO\_PI} \\
&V_X = \text{II\_CP} \hspace{1cm}& 
&X = \text{II\_Q} \hspace{1cm}& 
&P_X = \text{II\_PI} \\
&V_Y = \text{VA\_CP} \hspace{1cm}& 
&Y = \text{VA\_Q} \hspace{1cm}& 
&P_Y = \text{VA\_PI} \\
&V_L = \text{LAB} \hspace{1cm}& 
&L = \text{LAB\_QI} \hspace{1cm}& 
&K = \text{CAP\_QI}. \\
\end{align*}
Note that the yearly industry-level TFP growth available in the KLEMS data VAConTFP, does not correspond to Eq. \eqref{eq_tfp_translog} but is defined as 
\[
\text{LP1ConTFP}_{ic,t} =
     \hat{Q}_{ic,t} -
     \hat{L}_{ic,t} \cdot \alpha_{L,ic,t-1} -
    \hat{K}_{ic,t} \cdot \alpha_{K,ic,t-1}-
     \hat{X}_{ic,t} \cdot \alpha_{X,ic,t-1}.
\]
We recompute TFP using the more precise equation (Eq. \eqref{eq_tfp_translog}).

%Note that the long term TFP log growth rate is not perfectly equal to the eq \ref{eq_aihat_go} with long term growth rate. \ER{To be checked why}

\subsubsection{STAN} The \textit{STAN} database is from the OECD, it provides data for 23 countries covering 130 sectors.
\ER{Add outliers}

\begin{align*}
&V_Q=\text{PROD} \hspace{1cm}&
&Q=\text{PRDK} \hspace{1cm}&
&P_Q=\text{PRDP}\\
&V_X=\text{INTI} \hspace{1cm}&
&X=\text{INTK} \hspace{1cm}&
&P_X=\text{INTP}\\
&V_Y=\text{VALU} \hspace{1cm}&
&Y=\text{VALK} \hspace{1cm}&
&P_Y=\text{VALP}\\
&V_L=\text{LABR} \hspace{1cm}&
&L=\text{HRSN} \hspace{1cm}&
&K=\text{CPNK}\\
\end{align*}

When hours worked (HRSN) are missing, we substitute them with the employment variable (EMPN). To avoid inconsistent switching between the two measures, this replacement is applied only when the entire time series of HRSN is unavailable for a given industry. We use on the value added price index (VALP) of the aggregate industry “D01T99” as a proxy for country-specific GDP deflator.

\paragraph{Consistency check} To assess consistency across data sources, we compare the cumulative log growth of capital input, labour input, gross output, and total factor productivity (TFP) between KLEMS and STAN for matching industries over periods where both datasets overlap. The results show a high degree of alignment, suggesting that the measures are broadly consistent across the two sources.


\subsubsection{NBER CES} 

We use the NBER-CES dataset, following the 2012 NAICS classification at the 6-digit level. The dataset provides 61 years of data (1958–2018) for 364 industries within the U.S. manufacturing sector. We take total factor productivity (TFP) directly from the series \text{tfp5}, which corresponds to TFP levels constructed from logged growth rates using five inputs: production worker hours, non-production worker hours, capital, energy, and materials \citep{beckerNBERCESManufacturingIndustry2013}. We do not recompute TFP which is computed according to Eq. \eqref{eq_tfp_translog} \citep{beckerNBERCESManufacturingIndustry2013}.

\begin{align*}
A_{i,t,1} &= \text{tfp5}, \\
V_Q &= \text{vship}, \\
P_Q &= \text{piship}.
\end{align*}

\noindent
We compute real output as the ratio of nominal shipments to the corresponding price index,\footnote{See Section 4 of \cite{beckerNBERCESManufacturingIndustry2013}.}
\begin{equation}
    Q = \frac{\text{vship}}{\text{piship}} = \frac{V_Q}{P_Q}.
\end{equation}

\noindent
To deflate nominal values, we use the GDP deflator published by the BEA (series code \text{A191RG}).

\medskip

\noindent
We exclude several outlier sectors that exhibit implausible price spikes—specifically, sectors 334111, 334112, 334118, 334413, and 334613 (within computer manufacturing), and sectors 325312 and 325314 (“Phosphatic Fertilizer Manufacturing” and “Fertilizer (Mixing Only) Manufacturing”, respectively)—all of which display a one-time price increase exceeding 300\% in 2004.




% Bibliography ---------------------------

\tiny
\renewcommand\bibname{References}
\bibliographystyle{apalike}
\bibliography{z-bib}
\normalsize
\end{document}